% !PRINTED 217
\weberpage{217}
%<<CONT
gleichung für einen kleineren Divisor ist. Im Gegensatz hierzu
nennen wir die Gleichung, deren Wurzeln die \gesperrt{sämtichen} $x_{\mu,\mu'}$ % [sic]
\gesperrt{sind, die allgemeine Teilungsgleichung} für den Divisor $n$.

Durch irgend eine Substitution der Gruppe $\mathfrak{A}$
\[
S = \begin{pmatrix} a, & b\\ c, & \partial \end{pmatrix}
\]
geht $(1,0)$, $(0,1)$ in $(\partial, -c)$, $(-b, a)$ über, und wenn also $S$ nicht
die identische Substitution ist, so wird gewiß wenigstens eine der
beiden Wurzeln $x_{1,0}$, $x_{0,1}$ durch $S$ verändert. Daraus ergibt sich,
\gesperrt{daß die Galoissche Gruppe der eigentlichen Teilungsgleichung
genau dieselbe ist, wie die der allgemeinen,}
nur angewandt auf den Fall, daß $\mu, \mu', n$ keinen gemeinsamen
Teiler haben, nämlich, je nachdem die $n$ten Einheitswurzeln
adjungiert sind oder nicht, $\mathfrak{B}$ oder $\mathfrak{A}$.

Daraus folgt noch, \gesperrt{daß die eigentliche Teilungsgleichung
irreduzibel ist, selbst nach Adjunktion beliebiger Konstanten.}

Denn sind $\gamma, \delta$ irgend zwei Zahlen ohne gemeinsamen Teiler
mit $n$, so kann man $\alpha, \beta$ der Kongruenz
\[
\alpha\delta - \beta\gamma \equiv 1 \pmod{n}
\]
gemäß bestimmen und die in $\mathfrak{B}$ enthaltene Substitution
\[
T = \begin{pmatrix} \alpha, & \beta\\ \gamma, & \delta \end{pmatrix}
\]
auf jede rationale Gleichung zwischen den Wurzeln $x_{\mu,\mu'}$ anwenden.
Wenn also $x_{1,0}$ einer rationalen Gleichung (mit beliebigen konstanten
Koeffizienten)
\[
\Phi(x_{1,0}) = 0
\]
genügt, so genügt derselben Gleichung jede andere Wurzel $x_{\delta,-\gamma}$
der eigentlichen Teilungsgleichung, woraus die Irreduzibilität der
letzteren folgt.

\begin{center}
\S\,65. \gesperrt{Zurückführung der Teilungsgleichung}\\
\gesperrt{auf Transformationsgleichungen.}
\end{center}
Die Wurzeln der eigentlichen Teilungsgleichung lassen sich
in folgender Weise in \gesperrt{Reihen} anordnen. Man wähle nach Belieben
eine der Wurzeln:
\[
x_{\mu_1,\mu'_1} = \mathrm{sn}\left(\frac{4\mu_1 K + 4\mu'_1 iK'}{n}\right) = \mathrm{sn}\,\Omega_1.
\]
